\documentclass[12pt,a4paper]{article}

\usepackage[utf8]{inputenc}
\usepackage[brazilian]{babel}
\usepackage{setspace}
\usepackage{mathptmx}
\usepackage{geometry}
\usepackage{url}
\usepackage{xcolor}
\usepackage{fancyhdr}
\usepackage{longtable}
\usepackage{enumitem}


\geometry{a4paper, left=2cm, right=2cm, bottom=4cm, top=3cm}
\setstretch{1.5}
\setlength{\parindent}{0pt}

\pagestyle{fancy}
\fancyhf{}

% ------------------------------------------------
% HEADER
% ------------------------------------------------
\lhead{Ata de Reunião}
\chead{\meetingcode}
\rhead{\meetingdate}

% ------------------------------------------------
% FOOTER
% ------------------------------------------------
\fancyfoot[L]{Confidencial}
\fancyfoot[C]{$\copyright$ \companyname, \the\year}
\fancyfoot[R]{Página \thepage}

\renewcommand{\footrulewidth}{0pt}

% ------------------------------------------------
% MACROS DA REUNIÃO
% ------------------------------------------------

\newcommand{\meetingtitle}{}
\newcommand{\meetingtype}{}
\newcommand{\meetingdate}{}
\newcommand{\meetingtime}{}
\newcommand{\meetinglocation}{}
\newcommand{\facilitator}{}
\newcommand{\notetaker}{}
\newcommand{\participants}{}
\newcommand{\meetingobjective}{}
\newcommand{\meetingcode}{}
\newcommand{\companyname}{}

\newcommand{\setmeetingtitle}[1]{\renewcommand{\meetingtitle}{#1}}
\newcommand{\setmeetingtype}[1]{\renewcommand{\meetingtype}{#1}}
\newcommand{\setmeetingdate}[1]{\renewcommand{\meetingdate}{#1}}
\newcommand{\setmeetingtime}[1]{\renewcommand{\meetingtime}{#1}}
\newcommand{\setmeetinglocation}[1]{\renewcommand{\meetinglocation}{#1}}
\newcommand{\setfacilitator}[1]{\renewcommand{\facilitator}{#1}}
\newcommand{\setnotetaker}[1]{\renewcommand{\notetaker}{#1}}
\newcommand{\setparticipants}[1]{\renewcommand{\participants}{#1}}
\newcommand{\setobjective}[1]{\renewcommand{\meetingobjective}{#1}}
\newcommand{\setmeetingcode}[1]{\renewcommand{\meetingcode}{#1}}
\newcommand{\setcompanyname}[1]{\renewcommand{\companyname}{#1}}


% ------------------------------------------------
% INÍCIO DO DOCUMENTO
% ------------------------------------------------

\begin{document}

\setmeetingcode{MIIA-ATA-\textcolor{blue}{HASH (gerar hash correto)}}

\setmeetingtitle{\textcolor{blue}{Reunião de Planejamento da Sprint}}
\setmeetingtype{\textcolor{blue}{SCRUM}}
\setmeetingdate{\textcolor{blue}{04/03/2026}}
\setmeetingtime{\textcolor{blue}{14:00 -- 15:00}}
\setmeetinglocation{\textcolor{blue}{Google Meet}}
\setfacilitator{\textcolor{blue}{João Silva - nome de quem conduziu a reunião}}
\setnotetaker{\textcolor{blue}{André Morielo - nome de quem redigiu a ata}}
\setcompanyname{Instituto de Estudos Avançados (IEAv)}

\setparticipants{
\begin{itemize}
\item \textcolor{blue}{João Silva}
\item \textcolor{blue}{André Morielo}
\item \textcolor{blue}{Maria Santos}
\end{itemize}
}

\setobjective{\textcolor{blue}{Resumir em uma a duas frases os objetivos da reunião. Idealmente não deve ser extenso.}}

% ------------------------------------------------
% TÍTULO
% ------------------------------------------------

\begin{center}
{\Large \textbf{\meetingtitle}}\\
\vspace{0.3em}
{\small Código da Ata: \meetingcode}
\end{center}

\vspace{1em}

\textbf{Tipo de Reunião:} \meetingtype

\textbf{Data:} \meetingdate

\textbf{Horário:} \meetingtime

\textbf{Local / Plataforma:} \meetinglocation

\textbf{Facilitador:} \facilitator

\textbf{Responsável pela Ata:} \notetaker

\vspace{1em}

\textbf{Participantes}

\participants

\vspace{1em}

\textbf{Objetivo da Reunião}

\meetingobjective

\vspace{1em}

\textbf{Pauta}

\begin{enumerate}
\item \textcolor{blue}{Tópico 1}
\item \textcolor{blue}{Tópico 2}
\item \textcolor{blue}{Tópico 3}
\end{enumerate}

\vspace{1em}

\textbf{Discussões}

\begin{itemize}
\item \textcolor{blue}{Parágrafo mais extenso se necessário, tratando de eventuais discussões da reunião. Podem mencionar especificamente quando um determinado participante se manifestou, sobre qual tópica a discussão se tratou, entre outros assuntos.}
\item \textcolor{blue}{Adicionar quantos items forem necessários.}
\end{itemize}

\vspace{1em}

\textbf{Decisões Tomadas}

\begin{itemize}
\item \textcolor{blue}{Aqui são decisões, e não ações. Exemplo: ``Foi decidido que os cards da Sprint eram adequados, e poderiam ser efetivamente criados nos boards de controle de tarefas''}
\item \textcolor{blue}{Outro exemplo de decisão: ``Decidiu-se que é necessário estudar mais a fundo o tópico de Behavior Trees''}
\end{itemize}

\vspace{1em}

\textbf{Ações Definidas}

\begin{longtable}{|p{6cm}|p{4cm}|p{3cm}|}
\hline
\textbf{Ação} & \textbf{Responsável} & \textbf{Prazo} \\
\hline
 \textcolor{blue}{Criar os cards nos boards de controle de tarefas} & \textcolor{blue}{André Morielo} & \textcolor{blue}{01/04/2026} \\
\hline
 \textcolor{blue}{Estudar Behavior Trees} & \textcolor{blue}{João Silva} & \textcolor{blue}{A definir} \\
\hline
\end{longtable}

\vspace{1em}

\textbf{Bloqueios / Riscos}

\begin{itemize}
\item \textcolor{blue}{Geralmente associados às decisões ou ações. Exemplo: ``André Morielo sairá de férias no período, podendo impactar no prazo das atividades''}
\item \textcolor{blue}{Podem ser físicos também. Exemplo: ``Não há computador disponível para executar Behaviro Trees, o que pode impactar o estudo adequado dessas estruturar''}
\end{itemize}

\vspace{2em}

\textbf{Encerramento}

Nada mais havendo a tratar, a reunião foi encerrada.

\vspace{3em}

\textbf{\textcolor{blue}{Assinaturas - a assinatura de todos é opcional}}

\vspace{1em}

\begin{longtable}{|p{6cm}|p{9cm}|}
\hline
\textbf{Nome} & \textbf{Assinatura} \\
\hline
João Silva & \\[2.5cm]
\hline
André Morielo & \\[2.5cm]
\hline
Maria Santos & \\[2.5cm]
\hline
\end{longtable}

\end{document}